
% =========================================================
\section{Local structure and clipping-independent stability}
\label{sec:local-structure}
% =========================================================

We first establish a structural fact used throughout the paper: for
every fixed $c>0$, clipping is completely inactive in a sufficiently
small neighborhood of $x_*$, so the clipped and unclipped gradient
maps have identical local germs at $x_*$.  Consequently, clipping
cannot alter the linear stability threshold, the local invariant
structures, or any smooth local normal-form coefficient at $x_*$;
its influence can appear only when an orbit reaches a
finite-amplitude region intersecting the switching hypersurface
$\Sigma_c$.

% ---------------------------------------------------------
\subsection{Fixed points}
\label{subsec:fixed-points}
% ---------------------------------------------------------

We begin with a global observation.

\begin{proposition}[Fixed-point invariance]
\label{prop:fixed-point-invariance}
Let $\eta>0$ and $c>0$. Then
\begin{equation}
\operatorname{Fix}(F_{\eta,c})
=
\operatorname{Crit}(V)
:=
\left\{
x\in\mathcal U:
\nabla V(x)=0
\right\}.
\label{eq:fixed-critical-equivalence}
\end{equation}
In particular, the set of fixed points is independent of both
$\eta$ and $c$.
\end{proposition}

\begin{proof}
If $\nabla V(x)=0$, then
\[
\mathcal C_c^\diamond(\nabla V(x))=0,
\]
and hence
\[
F_{\eta,c}(x)
=
x-\eta\mathcal C_c^\diamond(\nabla V(x))
=
x.
\]
Thus every critical point of $V$ is a fixed point of $F_{\eta,c}$.

Conversely, suppose that
\[
F_{\eta,c}(x)=x.
\]
Since $\eta>0$, we obtain
\[
\mathcal C_c^\diamond(\nabla V(x))=0.
\]
By the definition of the radial clipping operator
\eqref{eq:clipping-operator},
\[
\mathcal C_c^\diamond(g)=0
\quad\Longleftrightarrow\quad
g=0.
\]
Therefore
\[
\nabla V(x)=0.
\]
This proves \eqref{eq:fixed-critical-equivalence}.
\end{proof}

\begin{remark}
\label{rem:no-spurious-fixed-points}
Proposition~\ref{prop:fixed-point-invariance} shows that norm clipping
does not create spurious equilibria.  All dynamical effects of
clipping therefore arise from the behavior of nonstationary
trajectories rather than from a modification of the critical-point
set of the potential.
\end{remark}

% ---------------------------------------------------------
\subsection{Quantitative inactivity of clipping}
\label{subsec:local-inactivity}
% ---------------------------------------------------------

The next lemma gives a quantitative local estimate.

For a linear map $A:\mathbb R^d\to\mathbb R^d$, define the mixed
operator norm
\begin{equation}
\|A\|_{2\to\diamond}
:=
\sup_{\|z\|_2=1}
\|Az\|_\diamond.
\label{eq:mixed-operator-norm}
\end{equation}

\begin{lemma}[Local gradient estimate]
\label{lem:local-gradient-estimate}
Under \textbf{(H1)--(H2)}, there exist $r_0>0$ and $L_*>0$ such that
\[
\overline{B_{r_0}(x_*)}
\subset\mathcal U
\]
and
\begin{equation}
\|\nabla V(x)\|_\diamond
\leq
L_*\|x-x_*\|_2
\qquad
\text{for all }
x\in B_{r_0}(x_*).
\label{eq:local-gradient-bound}
\end{equation}
\end{lemma}

\begin{proof}
Since $\mathcal U$ is open and $x_*\in\mathcal U$, choose
$r_0>0$ sufficiently small so that
\[
\overline{B_{r_0}(x_*)}
\subset\mathcal U.
\]
By \textbf{(H1)}, the Hessian $D^2V$ is continuous. Hence
\begin{equation}
L_*
:=
\sup_{y\in\overline{B_{r_0}(x_*)}}
\|D^2V(y)\|_{2\to\diamond}
<\infty.
\label{eq:L-star-definition}
\end{equation}

For $x\in B_{r_0}(x_*)$, the line segment joining $x_*$ to $x$
lies in $B_{r_0}(x_*)$. Since $\nabla V(x_*)=0$, the fundamental
theorem of calculus gives
\begin{align}
\nabla V(x)
&=
\nabla V(x)-\nabla V(x_*)
\nonumber\\
&=
\int_0^1
D^2V\bigl(x_*+t(x-x_*)\bigr)(x-x_*)\,dt.
\label{eq:gradient-integral}
\end{align}
Therefore
\begin{align}
\|\nabla V(x)\|_\diamond
&\leq
\int_0^1
\left\|
D^2V\bigl(x_*+t(x-x_*)\bigr)
\right\|_{2\to\diamond}
\|x-x_*\|_2\,dt
\nonumber\\
&\leq
L_*\|x-x_*\|_2.
\end{align}
This proves \eqref{eq:local-gradient-bound}.
\end{proof}

For every fixed $c>0$, define
\begin{equation}
r_c
:=
\min
\left\{
r_0,
\frac{c}{2L_*}
\right\}.
\label{eq:rc-definition}
\end{equation}

\begin{corollary}[Clipping inactivity]
\label{cor:clipping-inactivity}
For every $c>0$,
\begin{equation}
B_{r_c}(x_*)
\subset\Omega_c^-.
\label{eq:ball-inside-smooth-region}
\end{equation}
Equivalently,
\begin{equation}
B_{r_c}(x_*)\cap\Sigma_c
=
\varnothing.
\label{eq:no-switching-near-critical}
\end{equation}
\end{corollary}

\begin{proof}
If $x\in B_{r_c}(x_*)$, then by
Lemma~\ref{lem:local-gradient-estimate},
\[
\|\nabla V(x)\|_\diamond
\leq
L_*\|x-x_*\|_2
<
L_*r_c
\leq
\frac{c}{2}
<
c.
\]
Hence $x\in\Omega_c^-$.
\end{proof}

\begin{remark}[Dependence on the clipping threshold]
\label{rem:shrinking-neighborhood}
The neighborhood in Corollary~\ref{cor:clipping-inactivity} is not
uniform as $c\to0^+$.  The explicit choice
\eqref{eq:rc-definition} has
\[
r_c=O(c).
\]
This shrinking local scale is essential for the later
smooth-to-switching analysis: although clipping is invisible in a
neighborhood of $x_*$ for every fixed $c>0$, the boundary of that
neighborhood approaches the critical point as $c\to0^+$.
\end{remark}

% ---------------------------------------------------------
\subsection{Local identity of the clipped and smooth maps}
\label{subsec:local-identity}
% ---------------------------------------------------------

We can now state the basic structural theorem.

\begin{theorem}[Local identity theorem]
\label{thm:local-identity}
Assume \textbf{(H1)--(H2)}. For every fixed $c>0$, there exists
$r_c>0$ such that, for every $\eta>0$,
\begin{equation}
F_{\eta,c}(x)
=
G_\eta(x)
=
x-\eta\nabla V(x)
\qquad
\text{for all }
x\in B_{r_c}(x_*).
\label{eq:local-map-identity}
\end{equation}
The radius $r_c$ may be chosen independently of $\eta$.
\end{theorem}

\begin{proof}
By Corollary~\ref{cor:clipping-inactivity},
\[
B_{r_c}(x_*)
\subset\Omega_c^-.
\]
Hence, for every $x\in B_{r_c}(x_*)$,
\[
\|\nabla V(x)\|_\diamond<c.
\]
By definition of the clipping operator,
\[
\mathcal C_c^\diamond(\nabla V(x))
=
\nabla V(x).
\]
Substitution into \eqref{eq:clipped-gradient-map} yields
\[
F_{\eta,c}(x)
=
x-\eta\nabla V(x)
=
G_\eta(x).
\]
The construction of $r_c$ in \eqref{eq:rc-definition} does not
involve $\eta$, so the same neighborhood works for every $\eta>0$.
\end{proof}

Theorem~\ref{thm:local-identity} is stronger than equality of the
linearizations: the two maps agree as functions on an open
neighborhood of $x_*$.

% ---------------------------------------------------------
\subsection{Jet invariance and local normal forms}
\label{subsec:jet-invariance}
% ---------------------------------------------------------

The local identity immediately implies equality of all derivatives
allowed by the regularity of $V$.

\begin{corollary}[Local jet and normal-form invariance]
\label{cor:jet-invariance}
\label{cor:flip-coefficient-independence}
Under \textbf{(H1)--(H4)}, for every fixed $c>0$ and $\eta>0$:
\begin{enumerate}[label=\textnormal{(\roman*)}]
\item
every local Taylor coefficient of $F_{\eta,c}$ and $G_\eta$ at
$x_*$ that is defined under the assumed regularity is identical;
in particular,
\begin{equation}
DF_{\eta,c}(x_*)
=
I-\eta H_*;
\label{eq:clipped-linearization}
\end{equation}
\item
the local center-manifold equations and all smooth flip
normal-form coefficients of $F_{\eta,c}$ at $x_*$ are identical to
those of $G_\eta$; in particular, the coefficient $\ell_*$
defined in \eqref{eq:ell-star} is independent of the clipping
threshold $c$.
\end{enumerate}
\end{corollary}

\begin{proof}
The two maps coincide on a neighborhood of $x_*$ by
Theorem~\ref{thm:local-identity}, so all derivatives at $x_*$ that
are defined coincide.  Smooth normal-form coefficients depend only
on a finite jet at the bifurcation point; hence the center-manifold
invariance equations and reduced normal forms are identical.
\end{proof}

\begin{remark}
\label{rem:local-versus-finite-amplitude}
Corollary~\ref{cor:flip-coefficient-independence} is a central
structural distinction in the present problem.  Clipping cannot
modify the birth mechanism of a sufficiently small smooth
period-two orbit.  It can affect that orbit only after its amplitude
has grown sufficiently for one of its points to reach $\Sigma_c$.
The latter is a finite-amplitude event and will be analyzed in
Sections~\ref{sec:switching-geometry} and
\ref{sec:separation-law}.
\end{remark}

% ---------------------------------------------------------
\subsection{Exact local stability threshold}
\label{subsec:local-stability-threshold}
% ---------------------------------------------------------

We next determine the stability of $x_*$.  The threshold
$\eta_f=2/\lambda_*$ below is the classical maximal step size of
(unclipped) gradient descent \cite{Polyak1987,Nesterov2018};
Theorem~\ref{thm:local-stability-threshold} shows that it is
preserved exactly by clipping.

\begin{theorem}[Clipping-independent local stability threshold]
\label{thm:local-stability-threshold}
Assume \textbf{(H1)--(H3)} and let
\[
\eta_f=\frac{2}{\lambda_*}.
\]
Then, for every fixed $c>0$, the following statements hold.

\begin{enumerate}[label=\textnormal{(\roman*)}]

\item
If
\begin{equation}
0<\eta<\eta_f,
\label{eq:stable-eta-range}
\end{equation}
then $x_*$ is a locally asymptotically stable fixed point of
$F_{\eta,c}$.

\item
If
\begin{equation}
\eta>\eta_f,
\label{eq:unstable-eta-range}
\end{equation}
then $x_*$ is unstable for $F_{\eta,c}$.

\item
At
\begin{equation}
\eta=\eta_f,
\end{equation}
the fixed point is nonhyperbolic.  Its unique multiplier on the unit
circle is the simple multiplier $-1$ associated with $v_*$.

\end{enumerate}

Consequently, the first local stability threshold is exactly
\begin{equation}
\boxed{
\eta_f
=
\frac{2}{\lambda_{\max}(H_*)}
}
\label{eq:boxed-stability-threshold}
\end{equation}
and is independent of the clipping threshold $c$.
\end{theorem}

\begin{proof}
By Corollary~\ref{cor:jet-invariance},
\[
DF_{\eta,c}(x_*)
=
I-\eta H_*.
\]
Since $H_*$ is symmetric positive definite, the multipliers are
\[
\rho_i(\eta)
=
1-\eta\lambda_i,
\qquad
i=1,\ldots,d.
\]

Suppose first that
\[
0<\eta<\frac{2}{\lambda_*}.
\]
Since $0<\lambda_i\leq\lambda_*$,
\[
0<\eta\lambda_i<2,
\]
and hence
\[
-1
<
1-\eta\lambda_i
<
1
\qquad
\text{for every }i.
\]
Thus
\[
\max_i|\rho_i(\eta)|<1.
\]
The fixed point is therefore locally asymptotically stable by the
standard linearization theorem for discrete dynamical systems.

Now suppose that
\[
\eta>\frac{2}{\lambda_*}.
\]
Then
\[
\rho_d(\eta)
=
1-\eta\lambda_*
<
-1,
\]
so
\[
|\rho_d(\eta)|>1.
\]
Hence $x_*$ is unstable.

Finally, if
\[
\eta=\eta_f=\frac{2}{\lambda_*},
\]
then
\[
\rho_d(\eta_f)=-1.
\]
By \textbf{(H3)}, for every $i<d$,
\[
0<\lambda_i<\lambda_*,
\]
and therefore
\[
-1
<
1-\frac{2\lambda_i}{\lambda_*}
<
1.
\]
Thus $-1$ is the unique multiplier on the unit circle, and it is
simple.
\end{proof}

\begin{remark}
\label{rem:boundary-inconclusive-linearly}
Theorem~\ref{thm:local-stability-threshold} deliberately makes no
stability claim at $\eta=\eta_f$.  Linearization is inconclusive
when the critical multiplier has modulus one.  The nonlinear
behavior at this parameter is governed by the flip coefficient
$\ell_*$ and will be analyzed in Section~\ref{sec:flip-bifurcation}.
\end{remark}

% ---------------------------------------------------------
\subsection{Local invariant structures}
\label{subsec:local-invariant-structures}
% ---------------------------------------------------------

The local identity also transfers invariant structures from the
smooth gradient map to the clipped map.  For a map $T$ and a radius
$r>0$, define the local stable set relative to $B_r(x_*)$ by
\begin{equation}
W^s_r(x_*;T)
:=
\left\{
x\in B_r(x_*):
T^n(x)\in B_r(x_*)\ \forall n\geq0,
\quad
T^n(x)\to x_*
\right\}.
\label{eq:relative-local-stable-set}
\end{equation}

\begin{proposition}[Equality of local stable dynamics]
\label{prop:local-stable-dynamics}
Let $c>0$, and choose $0<r\leq r_c$.
Then, for every $\eta>0$,
\begin{equation}
W^s_r(x_*;F_{\eta,c})
=
W^s_r(x_*;G_\eta).
\label{eq:stable-set-equality}
\end{equation}
\end{proposition}

\begin{proof}
The two maps coincide on $B_r(x_*)$ by
Theorem~\ref{thm:local-identity}, so the defining conditions in
\eqref{eq:relative-local-stable-set} are equivalent for the two
maps.
\end{proof}

When $x_*$ is hyperbolic, this equality gives equality of the
corresponding stable-manifold germs.  If, in addition,
\begin{equation}
\det(I-\eta H_*)\neq0,
\label{eq:local-diffeomorphism-condition}
\end{equation}
then both maps are local diffeomorphisms near $x_*$, and the same
argument gives equality of their local unstable-manifold germs.

At the critical parameter $\eta=\eta_f$, the fixed point is
nonhyperbolic. Nevertheless, local identity still implies equality
of the corresponding center dynamics.

\begin{proposition}[Center-manifold invariance at the flip threshold]
\label{prop:center-manifold-invariance}
Assume \textbf{(H1)--(H3)} and set $\eta=\eta_f$.
For every fixed $c>0$, a sufficiently small local center manifold of
$G_{\eta_f}$ through $x_*$ is also a local center manifold of
$F_{\eta_f,c}$.

Moreover, the reduced maps on this common center manifold are
identical.
\end{proposition}

\begin{proof}
By Theorem~\ref{thm:local-identity}, the maps
$F_{\eta_f,c}$ and $G_{\eta_f}$ coincide in a neighborhood of
$x_*$. Choose a local center manifold of $G_{\eta_f}$ contained in
that neighborhood. Its local invariance equation is therefore also
satisfied by $F_{\eta_f,c}$. Since the restrictions of the two maps
to the manifold coincide pointwise, the corresponding reduced
one-dimensional dynamics are identical.
\end{proof}

\begin{remark}[Nonuniqueness of center manifolds]
\label{rem:center-manifold-nonuniqueness}
Center manifolds need not be unique as sets; the content of
Proposition~\ref{prop:center-manifold-invariance} is only that the
two maps admit one common choice with identical reduced dynamics.
\end{remark}

% ---------------------------------------------------------
\subsection{Consequences for the smooth-to-switching mechanism}
\label{subsec:local-structure-consequences}
% ---------------------------------------------------------

\begin{corollary}[Local bifurcation structure is clipping-independent]
\label{cor:local-bifurcation-clipping-independent}
Under \textbf{(H1)--(H4)}, for every fixed $c>0$:

\begin{enumerate}[label=\textnormal{(\roman*)}]

\item
the fixed point $x_*$ is unchanged by clipping;

\item
its first local stability threshold is
\[
\eta_f=\frac{2}{\lambda_*};
\]

\item
the critical eigendirection is $v_*$;

\item
the critical multiplier crosses $-1$ transversally at $\eta_f$;

\item
the local center-manifold reduction is identical to that of the
unclipped gradient map;

\item
the supercritical flip coefficient $\ell_*$ is independent of $c$.

\end{enumerate}
\end{corollary}

\begin{proof}
Items (i)--(vi) follow respectively from
Proposition~\ref{prop:fixed-point-invariance},
Theorem~\ref{thm:local-stability-threshold},
\eqref{eq:critical-eigenpair},
\eqref{eq:transversal-crossing},
Proposition~\ref{prop:center-manifold-invariance}, and
Corollary~\ref{cor:flip-coefficient-independence}.
\end{proof}

The key implication is that $c$ does not enter the local birth
mechanism of the period-doubled branch: its first dynamical effect
must occur at a distinct finite-amplitude event, when the
bifurcating orbit reaches the endogenous switching hypersurface
$\Sigma_c$.

